% ===== §3 Synthesis: the axiom of sustainability — geometric phase =====
\section{Synthesis: The Axiom of Sustainability and the Criterion of Geometric Phase}
\label{p09:sec:synthesis}

\noindent
The synthesis must supply what the cycle, by itself, withholds: a criterion that distinguishes the good circle from the vicious one without appeal to the bare fact of self-continuation, which both share. I shall argue that the criterion is \emph{geometric}---that the difference between the good and the vicious cycle is, with surprising exactness, the difference between a \emb{circle} and a \emb{spiral}, between a return that accrues nothing and a return that comes home bearing an irreducible remainder. And I shall argue that this is not a metaphor but a structural isomorphism, for which the physics and geometry of \emph{anholonomy}---of the geometric, or Berry--Pancharatnam, phase---supplies the precise form.

\subsection{Return to the root, and the difference between the circle and the spiral}
\label{p09:sec:returnroot}

Begin with the Daoist intuition the prelude already sounded. The \zh{反} of \zh{反者道之动} is not mere oscillation, a swinging back and forth between poles; it is \zh{返}---a \emph{returning}, a coming-home to the root. ``The myriad things in their teeming each return to their root'' (\zh{夫物芸芸，各复归其根}, \emph{Daodejing} ch.~16; \cite{laozi1963}). This returning must be sharply distinguished from a second figure with which it is easily confused: the monotone linear accumulation that Hegel named the \emph{bad infinite} (\emph{schlechte Unendlichkeit})---the mere ``and then, and then,'' the endless additive extrapolation that never comes home but only extends \parencite{hegel2010}. Capital's M--C--M$'$ is bad-infinite in just this sense: it does not return to its root; it extends without rest toward a more that is never enough.

So we already have two figures that must not be conflated: the \emph{returning} of the Dao (return-to-source) and the \emph{accumulating} of the bad infinite (linear extension). But this is not yet the distinction we need, for the antithesis showed that a returning cycle, too, can be vicious (the fair stasis that returns value yet merely repeats). We need a third figure, and it is here that geometry supplies what the two intuitions cannot.

\subsection{The geometric phase: return without return-to-the-same}
\label{p09:sec:geom-phase}

The decisive phenomenon is this. Consider a system carried around a \emph{closed loop} in some space of parameters or situations---carried, that is, so that at the end it returns to the very point from which it set out. One might expect that, having come back to the same point, the system is in the same state. For an important class of systems this expectation is false. When the dynamical (the trivially accumulated) phase is subtracted, the system does \emph{not} return to its initial state: it has accrued a residual---a \emb{geometric phase}---determined purely by the geometry of the path, by the ``area'' the loop encloses, which is to say by the integral of a curvature over the surface the loop bounds. This residual is the phenomenon of \emph{anholonomy}: \emb{the loop closes in the base space, and yet does not close in the fibre.}\footnote{The canonical physical instance is Berry's: a quantum system carried adiabatically around a closed loop in parameter space returns with a phase factor that is not dynamical but geometric---an observable, gauge-invariant quantity equal to the flux of a curvature through the enclosed surface \parencite{berry1984}. The classical optical antecedent is Pancharatnam's \parencite{pancharatnam1956}, for the phase of polarized light carried around a circuit on the Poincaré sphere. What matters for us is the structure, not the particular physics: a closed circuit in the base, an open one in the fibre, the gap between them fixed by an enclosed curvature.}

This is, with a precision that startled the author when it first appeared, the exact form of ``return to the root, yet not to the origin.'' The dialectical movement---thesis, antithesis, synthesis---goes once around and comes back to its word; ``the Buddha speaks of a world, which is then no world, and is thus named a world'' (\zh{如来说世界，则非世界，是名为世界}, the \emph{Diamond Sutra}): the third term returns to ``world,'' to the very word of the first, and yet, having passed through the negation of the negation, it is no longer the naive ``world'' of the outset---it carries a remainder, a sublimation, a geometric phase. \emb{The return happens in the base space (one comes home, to the same person, the same daily round); the sublimation happens in the fibre (one brings home what was not there at setting out).}

We can now state the criterion the synthesis was charged to supply.

\begin{claim}[The geometric criterion of the good and vicious cycle]
Let the state of a relation, or of a value-process, evolve on a fibre bundle: the \emph{base} the space of relational content or situation, the \emph{fibre} the phase of value, of recognition, of jouissance. A cycle of reproduction is \emph{vicious} when its holonomy is zero or negative---when, going once around, it either returns to exactly the point of origin (the stalled oscillation; Hegel's bad infinite is this zero-creating monotone) or suffers a net loss of phase (a net extraction each turn, the rose exhausted). A cycle is \emph{good} when its holonomy is positive---when, in the base, it returns to its root, while in the fibre it has accrued a positive geometric phase. \emb{The circle is the vicious cycle (zero-phase repetition); the spiral is the good cycle (positive holonomy).}
\end{claim}

Note what this criterion accomplishes that the bare notion of self-continuation could not. It \emb{frees ``self-as-end'' from the work of distinguishing good and ill}. Both capital's M--C--M$'$ and the good relational cycle ``take their own continuation as their end''; the difference is not in the autotelic form but in the \emph{mechanism of the phase}. Capital's increment is counterfeited by extraction from an outside (an open system masquerading as a closed loop, in the language of \xref{p09:sec:antithesis}); the good cycle's positive phase is \emph{endogenous}---produced by the curvature of the path itself, by the internal tension-structure of the relation, with no outside drawn upon. This is the geometric form of the distinction between the \emph{stolen} and the \emph{internal} phase that will govern the rest of the paper.

\subsection{The axiom of sustainability: the advance upon generative justice}
\label{p09:sec:axiom}

With the criterion in hand, the paper's central constructive claim---its advance upon Eglash---can be stated.

\begin{claim}[The axiom of sustainability]
Generative justice requires not only that value return to its creators (Eglash's spatial criterion) but that this return suffice, each turn, to reproduce the conditions of generation themselves---and, more strongly, that the cycle accrue a positive geometric phase, a sublimation, with each turn. A cycle that returns value to its creators and yet exhausts the very basis of its own reproduction is just in its distribution and unsustainable in its time. \emb{The axiom of sustainability is the temporal dimension that the spatial criterion lacks: the good circle must return to its root \emph{and} sublimate---it must be a spiral.} This is the paper's contribution to generative justice: not an application of it, but the addition to it of a criterion of time.
\end{claim}

The axiom illuminates, too, why \emph{wu wei}---non-action---will turn out (in \xref{p09:sec:praxis}) to be the fitting practical disposition. If the geometric structure of a system, the curvature of its bundle, is already such that its holonomy is positive, then \emb{the sublimation happens of itself}, and requires no external forcing. \emph{Wu wei} is not inaction but the refusal to extract, to force-close, to impose an external linear end upon a cycle that would, left to its own curvature, sublimate of itself. Forced action---the contrivance of capital, of instrumental reason---is precisely the disease that, in straining to control the phase and wring out an increment, destroys the very curvature by which the system would naturally have produced its sublimation. \emb{Let the cycle be the spiral it would of itself become, and it will sublimate.} This is the geometric--Daoist statement of ``let the circulation of value become a good, natural circle, and it will perpetuate itself of itself.''

\subsection{Two honest caveats}
\label{p09:sec:caveats}

Two debts must be acknowledged at once, for the geometric language is powerful enough to seduce, and the paper had rather name its limits than be caught out at them.

\begin{claim}[First caveat: the phase requires a real bearer]
A geometric phase is physics, and not metaphor, only because the quantity carried around the loop is real---a phase with the structure of a circle (an $S^1$, or a more general group), additive, gauge-invariant, in principle observable. If we speak of ``the phase of value,'' we are obliged to answer: the phase of \emph{what}? What is it that has a periodic, additive, gauge-invariant structure, such that a holonomy can be spoken of at all? Absent a real bearer, the geometric phase is borrowed eloquence and no more. The paper does not discharge this debt here; it incurs it openly, and pays it only in \xref{p09:sec:core}, where the bearer is fixed as the sublimation of \emph{real value}---the jouissance generated and refluxing in the interpretive cycle. Until then the geometric language is to be read as a structural analogy with a forward promise, not a completed formalism. The technical formalization, with its honest treatment of the bearer and the dimensionality, is deferred to the appendix and to the separate Value-Foam study it points to.
\end{claim}

\begin{claim}[Second caveat: ``sublimation $=$ positive phase'' must be argued, not stipulated]
It is not permitted to \emph{define} sublimation as positive holonomy and then announce an isomorphism; that would be to prove by naming. One must characterize, independently, a quantity of ``sublimation'' or ``surplus''---for instance, a margin of sustainability: the net reproduction, each cycle, of the conditions of generation, in excess of what is consumed---and then \emph{show} that, under a suitable bundle formulation, it corresponds to the sign of the holonomy. To weld the two by definition is circular (a vicious circle, in the most literal sense). This paper states the requirement and meets it only in part, in \xref{p09:sec:core}; it flags, here, that the correspondence is a claim to be earned and not a convenience to be assumed. The reader will find, in \xref{p09:sec:turn}, the paper calling itself to account on exactly this point.
\end{claim}

These caveats are not weaknesses to be hidden but the very form of intellectual honesty the paper's epistemology (\xref{p09:sec:epistemology}) will make thematic: that the geometric framework offers a structural \emph{characterization} of what the good circle \emph{is}, and not an operational \emph{procedure} for measuring it from without. To mistake the ``what it is'' for the ``how to measure'' is the positivist's temptation, and the next section is, in part, a guard against it.
